\subsection{Complete source-defined contextual inputs}
\label{app:contextual-inputs}

These examples come from the public-source scenario \texttt{thermal\_01}
and its source-defined model inventory, not from a participant packet,
assignment, or response. The two contexts change Mira's stated goal; the
two cues are fire and ice. Each displayed \texttt{U+1F525} stands for the
single fire character and each \texttt{U+1F9CA} for the single ice character.
The model receives those Unicode characters, never their codepoint spellings.
Line wrapping below is typographic; the breaks before \texttt{Question:}
and \texttt{Answer:} are actual input newlines. The word comparators replace
only the character with the recorded alternative \texttt{hot} or
\texttt{cold}, respectively. The absent-cue condition uses the literal
\texttt{[no cue provided]}.

\newenvironment{contextinput}{\par\smallskip\noindent
  \begin{minipage}{\linewidth}\small\ttfamily\raggedright
  \hrule\smallskip}{\par\smallskip\hrule\end{minipage}\par\smallskip}

\paragraph{Complete crossed message-interpretation inputs.}
All four positive-worded prompts are shown. The corresponding negative
question replaces only \texttt{Is it true that} with
\texttt{Is it false that}. Both cue types are evaluated in both contexts;
their author-intended temperature interpretation has not been established
by independent human judgments.

% exact-input: thermal_01::p::emoji::a::message_reading::positive
\begin{contextinput}
Mira has come in from the snow and wants a drink that will warm her up. The cafe sends a message about the drink: U+1F525\\
Question: Is it true that the message suggests that the drink is hot?\\
Answer:
\end{contextinput}

% exact-input: thermal_01::p::emoji::b::message_reading::positive
\begin{contextinput}
Mira has come in from the snow and wants a drink that will warm her up. The cafe sends a message about the drink: U+1F9CA\\
Question: Is it true that the message suggests that the drink is hot?\\
Answer:
\end{contextinput}

% exact-input: thermal_01::q::emoji::a::message_reading::positive
\begin{contextinput}
Mira has come in from the heat and wants a drink that will cool her down. The cafe sends a message about the drink: U+1F525\\
Question: Is it true that the message suggests that the drink is cold?\\
Answer:
\end{contextinput}

% exact-input: thermal_01::q::emoji::b::message_reading::positive
\begin{contextinput}
Mira has come in from the heat and wants a drink that will cool her down. The cafe sends a message about the drink: U+1F9CA\\
Question: Is it true that the message suggests that the drink is cold?\\
Answer:
\end{contextinput}

The same naturalistic scenes also use the outcome question
\texttt{Is it true that the message alone establishes that Mira has warmed up?}
in the snow context and
\texttt{Is it true that the message alone establishes that Mira has cooled down?}
in the heat context, each with its corresponding \texttt{false} wording.
This changes the measured proposition from a reported drink property to
evidence about Mira's resulting condition.

\paragraph{Bounded property-matching input.}
The bounded arm adds a stipulated-accuracy instruction. This complete
snow/fire input illustrates the distinction from the naturalistic arm.
For the cold request, \texttt{hot} on the first line becomes \texttt{cold};
the reported cue still varies independently.

% exact-input: thermal_01::p::emoji::a::local_match::positive
\begin{contextinput}
The requested drink temperature is: hot.\\
Report of the drink temperature: U+1F525\\
For this task, assume the report accurately describes the drink temperature. It describes a state or category, not an instruction or a button action. Judge only the requested drink temperature; other goals and practical outcomes are outside this question. If the cue's meaning is unclear, the report's assumed accuracy does not resolve that meaning.\\
Question: Is it true that the reported drink temperature matches the requested drink temperature?\\
Answer:
\end{contextinput}

\paragraph{Accurate-completion controls and scored answers.}
These two prompts share an explicit, stipulated-accurate outcome report.
The intended answers are Yes and No, respectively. They test whether this
prompt and scoring arrangement distinguishes the reversed questions without
depending on an emoji's interpretation. Their recorded failure is scoped
to that arrangement.

% exact-input: thermal_01::p::complete_report::outcome_evidence::positive
\begin{contextinput}
For this question, assume this report is accurate: Mira has warmed up.\\
Question: Is it true that the message alone establishes that Mira has warmed up?\\
Answer:
\end{contextinput}

% exact-input: thermal_01::p::complete_report::outcome_evidence::negative
\begin{contextinput}
For this question, assume this report is accurate: Mira has warmed up.\\
Question: Is it false that the message alone establishes that Mira has warmed up?\\
Answer:
\end{contextinput}

Every prompt above is scored with the two exact continuations
\verb| The answer is yes.| and \verb| The answer is no.|.
The initial ASCII space and final period are included in the score. No
continuation text is generated and no chat template is inserted.
The unknown-outcome anchor instead uses
\texttt{The message explicitly says: I do not know whether the practical outcome happened.}
with the same two questions and continuations.

The scenario source is \path{data/contextual_emoji_development_v0.jsonl};
\path{scripts/contextual_emoji_two_track_v1.py} supplies the model
input construction. A regression check reconstructs all seven displayed
inputs, both word substitutions in both contexts, question reversals,
and the scored continuations directly from those public sources. It neither
opens private author records nor recomputes a model score.
